% Part: lambda-calculus
% Chapter: lambda-definability
% Section: introduction

\documentclass[../../../include/open-logic-section]{subfiles}

\begin{document}

\olfileid{lam}{ldf}{int}
\olsection{परिचय}

पहिल्या नजरेला लॅम्डा कलन हे फलने दर्शवणाऱ्या
अभिव्यक्ती आणि त्यांचे इतर अभिव्यक्तींना उपयोजन
यांचे अतिशय अमूर्त कलन वाटते. या कलनाच्या
विन्यासात ही फलने कोणत्या विशिष्ट प्रकारच्या वस्तूंवर
असतात याचा निर्देश नाही; विशेषतः नैसर्गिक संख्यांचा
उल्लेखही नाही. त्याच्या मूलभूत क्रिया—उपयोजन आणि
लॅम्डा अमूर्तीकरण—सर्व प्रकारच्या फलनांवर लागू
होतात, केवळ नैसर्गिक संख्यांवरील फलनांवर नाही.

तरीही थोड्या कल्पकतेने लॅम्डा कलनात अंकगणितीय
फलने, म्हणजे नैसर्गिक संख्यांवरील फलने, परिभाषित
करता येतात. त्यासाठी प्रत्येक नैसर्गिक संख्या~$n \in
\Nat$ करिता विशेष $\lambd$-पद~$\num n$ परिभाषित
करू. ते~$n$ साठीचे \emph{चर्च अंक} आहे. चर्च अंकांना
अलाँझो चर्च यांचे नाव दिले आहे.

\begin{defn}
  $n \in \Nat$ असेल, तर संबंधित \emph{चर्च अंक}
  $\num{n}$ हे~$n$ दर्शवते:
  \[
    \num{n} \ident \lambd[fx][f^n(x)]
  \]
  येथे $f^n(x)$ म्हणजे $f$ हे~$x$ वर $n$ वेळा लागू
  केल्याचा परिणाम. उदाहरणार्थ, $\num{0}$ हे
  $\lambd[fx][x]$ आहे आणि $\num{3}$ हे
  $\lambd[fx][f(f(f\,x))]$ आहे.
\end{defn}

चर्च अंक $\num n$ हे दोन युक्तिवाद $f$ आणि $x$
स्वीकारून $f^n(x)$ परत करणारे फलन दर्शवणारे लॅम्डा
पद म्हणून संकेतित केलेले आहे. चर्च अंक स्पष्टपणे
सामान्य रूपात आहेत.

लॅम्डा कलनात नैसर्गिक संख्या दर्शवणे उपयोगी ठरते
तेव्हाच, जेव्हा त्यांच्यावर संगणनही करता येते.
चर्च अंकांवर लॅम्डा कलनात संगणन करणे म्हणजे
एखादे $\lambd$-पद~$F$ अशा चर्च अंकाला लागू करणे
आणि तयार झालेले पद~$F\, \num n$ सामान्य रूपात
न्यूनीत करणे. ते नेहमी सामान्य रूपात न्यूनीत झाले
आणि ते सामान्य रूप नेहमी एखादे चर्च अंक~$\num m$
असेल, तर संगणनाचा परिणाम~$m$ ही संख्या मानता येते.
मग~$F$ हे $f\colon \Nat \to \Nat$ फलन परिभाषित
करते असे मानता येते: $f(n) = m$ तेव्हाच, जेव्हा
$F\, \num n \red \num m$. चर्च–रॉसर गुणधर्मामुळे
सामान्य रूप अस्तित्वात असेल, तर ते एकमेव असते.
म्हणून $F\, \num n \red \num m$ असेल, तर
$F \, \num n$ चे दुसऱ्या कोणत्याही सामान्य रूपात,
विशेषतः दुसऱ्या चर्च अंकात, न्यूनीकरण होऊ शकत नाही.

उलट, $f\colon \Nat \to \Nat$ फलन दिले असता,
अशा रीतीने~$f$ परिभाषित करणारे पद $F$ आहे का,
असे विचारता येते. असे असेल, तर $F$ हे~$f$ ला
\emph{!!{lambda define}s} असे म्हणतो आणि $f$ हे
!!{lambda definable} आहे असे म्हणतो. ही संकल्पना
अनेकस्थानी आणि आंशिक फलनांपर्यंत व्यापक करता येते.

\begin{defn}
$f\colon \Nat^k \to \Nat$ असे मानू. सर्व
$n_0$, \dots,~$n_{k-1}$ साठी $f(n_0, \dots,
n_{k-1})$ परिभाषित असेल, तर
\[
F \, \num{n_0} \, \num{n_1} \dots \num{n_{k-1}} \red \num{f(n_0, n_1, \dots,
  n_{k-1})}
\]
आणि अन्यथा $F \, \num{n_0} \,
\num{n_1} \dots \num{n_{k-1}}$ ला सामान्य रूप
नसेल, तेव्हा लॅम्डा पद~$F$ हे $f$ ला
\emph{!!{lambda define}s} असे म्हणतो.
\end{defn}

स्थिर फलने हे एक अत्यंत सोपे उदाहरण आहे.
$C_k \ident \lambd[x][\num{k}]$ हे पद
$c_k\colon \Nat \to \Nat$ फलनाला
!!{lambda define}s; येथे $c_k(n) = k$.
कारण कोणत्याही~$n$ साठी $C_k \, \num n \ident
(\lambd[x][\num{k}])\num n \redone \num{k}$.
ओळख फलन $\lambd[x][x]$ ने !!{lambda defined} आहे.
अधिक गुंतागुंतीची फलने परिभाषित करणे अर्थातच
कठीण असते आणि त्यासाठी बरीच कल्पकता लागते.
त्यामुळे प्रत्येक संगणनक्षम फलन !!{lambda definable}
आहे हे काहीसे आश्चर्यकारक वाटू शकते. याचा उलट
दिशेचा निष्कर्षही खरा आहे: फलन !!{lambda definable}
असेल, तर ते संगणनक्षम असते.

\end{document}
